% =============================================================================
% Section 3 — Results
% =============================================================================
\section{Results}
\label{sec:results}

\subsection{Coupling parameter transition}
\label{sec:ka-transition}

Figure~\ref{fig:ka-transition} shows the Helmholtz number $ka$ and the
quadrupole coupling factor $(ka)^2$ as functions of frequency for the
canonical abdominal model ($R_\mathrm{eq} = \SI{0.157}{\metre}$).  At
the $n = 2$ flexural resonance (\SI{3.95}{\hertz}), $ka = 0.011$ and
$(ka)^2 = 1.3 \times 10^{-4}$.  By \SI{40}{\hertz}, $ka = 0.115$ and
$(ka)^2 = 1.3 \times 10^{-2}$---a hundredfold improvement.  At
\SI{80}{\hertz}, $ka = 0.23$ and $(ka)^2 = 5.3 \times 10^{-2}$.

\begin{table}[htbp]
  \centering
  \caption{Helmholtz number and coupling factor at selected frequencies.}
  \label{tab:ka-values}
  \begin{tabular}{@{}lllll@{}}
    \toprule
    $f$ [\si{\hertz}] & $ka$ & $(ka)^2$ &
      Ratio to $f_2$ & Regime \\
    \midrule
    3.95 (resonance) & 0.011 & $1.3 \times 10^{-4}$ & 1 & Deep Rayleigh \\
    20               & 0.057 & $3.3 \times 10^{-3}$ & 25 & Rayleigh \\
    40               & 0.115 & $1.3 \times 10^{-2}$ & 100 & Transition \\
    60               & 0.172 & $3.0 \times 10^{-2}$ & 230 & Transition \\
    80               & 0.230 & $5.3 \times 10^{-2}$ & 408 & Transition \\
    \bottomrule
  \end{tabular}
\end{table}

Despite this substantial improvement, the coupling factor remains well
below unity.  At \SI{80}{\hertz}, the effective driving pressure is
still only \SI{5.3}{\percent} of the incident pressure for $n = 2$.

\subsection{Off-resonance attenuation}
\label{sec:off-resonance}

The SDOF transfer function $H(r, \zeta)$ (equation~\ref{eq:sdof})
attenuates the response far from resonance.  Table~\ref{tab:sdof} shows
$H$ at representative frequencies.

\begin{table}[htbp]
  \centering
  \caption{SDOF transfer function at sub-bass frequencies
    ($f_2 = \SI{3.95}{\hertz}$, $\zeta = 0.125$).}
  \label{tab:sdof}
  \begin{tabular}{@{}llll@{}}
    \toprule
    $f$ [\si{\hertz}] & $r = f/f_2$ & $H(r, \zeta)$ & $H$ [\si{\decibel}] \\
    \midrule
    3.95 & 1.00 & 4.00 & +12.0 \\
    20   & 5.06 & 0.041 & $-27.8$ \\
    40   & 10.1 & 0.010 & $-40.0$ \\
    60   & 15.2 & 0.0044 & $-47.2$ \\
    80   & 20.3 & 0.0025 & $-52.2$ \\
    \bottomrule
  \end{tabular}
\end{table}

The off-resonance penalty dominates.  At \SI{40}{\hertz}, the
hundredfold improvement in coupling is offset by a hundredfold
reduction in transfer function, yielding net displacement comparable
to the resonance case.  At higher frequencies, the $1/r^2$ roll-off
wins convincingly.

\subsection{Tissue displacement at concert SPLs}
\label{sec:displacement-results}

Figure~\ref{fig:displacement-spectrum} presents tissue displacement as a
function of frequency for the three concert spectra.  Key results:

\begin{enumerate}
  \item \textbf{Electronic dance music} ($L_p \approx \SI{115}{\decibel}$
    at \SI{40}{\hertz}):  Peak displacement
    $\xi \approx \SI{e-4}{\micro\metre}$---four orders of magnitude below
    whole-body vibration perception thresholds.

  \item \textbf{Rock concert} ($L_p \approx \SI{108}{\decibel}$ at
    \SI{63}{\hertz}):  Peak displacement comparable to the EDM case.

  \item \textbf{Pipe organ} ($L_p \approx \SI{105}{\decibel}$ at
    \SI{40}{\hertz}):  Slightly lower displacement due to lower SPL.
\end{enumerate}

\begin{table}[htbp]
  \centering
  \caption{Tissue displacement at selected sub-bass frequencies for the
    EDM spectrum (peak genre).  Perception thresholds from ISO~2631-1.}
  \label{tab:displacement}
  \begin{tabular}{@{}llllll@{}}
    \toprule
    $f$ [\si{\hertz}] & $L_p$ [\si{\decibel}] & $ka$ &
      $\xi$ [\si{\micro\metre}] & $\xi_\mathrm{thresh}$
      [\si{\micro\metre}] & Ratio \\
    \midrule
    20 & 100 & 0.057 & --- & --- & --- \\
    31.5 & 112 & 0.091 & --- & --- & --- \\
    40 & 115 & 0.115 & --- & --- & --- \\
    50 & 114 & 0.144 & --- & --- & --- \\
    63 & 110 & 0.181 & --- & --- & --- \\
    80 & 106 & 0.230 & --- & --- & --- \\
    \bottomrule
  \end{tabular}
\end{table}

\noindent
\emph{Note: Numerical values in Table~\ref{tab:displacement} are
placeholders to be populated from the computational results.}

\subsection{Coupling transition band}
\label{sec:transition-band}

We define the ``coupling transition band'' as the frequency range in
which the ratio of airborne-induced displacement to perception threshold
exceeds \SI{1}{\percent}.  At \SI{115}{\decibel} SPL (the EDM peak),
no frequency in the \SIrange{1}{100}{\hertz} range meets this
criterion.  The peak ratio occurs near the $n = 2$ resonance
(\SI{3.95}{\hertz}), where the resonance amplification compensates
for the weaker coupling.

\subsection{Multi-modal contributions}
\label{sec:multimodal}

Higher-order flexural modes ($n = 3, 4, \ldots$) have resonance
frequencies closer to the sub-bass band but even weaker acoustic
coupling ($\propto (ka)^n$ with larger $n$).  Summing contributions from
$n = 2$ through $n = 6$ in quadrature increases the total displacement
by less than \SI{10}{\percent} at any frequency, confirming that the
$n = 2$ mode dominates throughout the sub-bass band.
